% Course book preamble, shared by every course book built from this kit.
% Editing this file restyles every course book at once.
% Rules for the content: ../../rules/explaining.md. Usage: README.md next to this file.
%
% Look: the internship report (Computer Modern, figure palette, black-outlined
% TikZ), plus a wide right margin that carries the definition of each difficult
% term next to the line where it first appears.

% Encoding and language
\usepackage[utf8]{inputenc}
\usepackage[T1]{fontenc}
\usepackage{lmodern}
\usepackage[english]{babel}
\usepackage{microtype}

% Page: 12.4 cm of text, 4.6 cm of margin notes on the right of every page
\usepackage{geometry}
\geometry{a4paper, left=2cm, textwidth=12.4cm, marginparsep=0.5cm,
  marginparwidth=4.6cm, top=2.5cm, bottom=2.5cm, headheight=14pt}
\usepackage{marginfix}          % margin notes never get lost or overlap
\usepackage{changepage}         % wide figures that run into the margin
\setlength{\parskip}{3pt}

% Mathematics
\usepackage{amsmath}
\usepackage{amssymb}

% Figures and tables
\usepackage{graphicx}
\graphicspath{{figures/}}
\usepackage{float}
\usepackage{caption}
\captionsetup{font=small,labelfont=bf,skip=6pt}
\usepackage{booktabs}
\usepackage{array}
\usepackage{tabularx}
\usepackage{multirow}
\usepackage{enumitem}
\setlist{itemsep=1pt,topsep=3pt}
\renewcommand{\arraystretch}{1.15}

% Diagrams and plots
\usepackage{xcolor}
\usepackage{tikz}
\usetikzlibrary{positioning,arrows.meta,calc,fit,backgrounds,shapes.geometric,decorations.pathreplacing}
\usepackage{pgfplots}
\pgfplotsset{compat=1.17}

% Figure palette of the internship report: each colour has its own luminance,
% so fills stay distinct in a greyscale print. No yellow.
\definecolor{figdark}{HTML}{1F4E79}
\definecolor{figblue}{HTML}{2E75B6}
\definecolor{figmid}{HTML}{2E8B8B}
\definecolor{figsoft}{HTML}{7FB3D5}
\definecolor{figlight}{HTML}{A9CCE3}
\definecolor{figpale}{HTML}{CFE2F0}
\definecolor{figbg}{HTML}{E6F0F8}
\definecolor{sevone}{HTML}{FCE9E3}
\definecolor{sevtwo}{HTML}{F3B3A0}
\definecolor{sevthree}{HTML}{D9604A}
\definecolor{sevfour}{HTML}{8E1B12}

% Shared TikZ styles: a box, a highlighted box, an arrow, a dashed group
\tikzset{
  bx/.style={draw=black, fill=figbg, rounded corners=1.5pt, align=center,
    font=\footnotesize, inner sep=4pt, minimum height=8mm},
  bxhi/.style={bx, fill=figpale},
  bxdark/.style={bx, fill=figdark, text=white},
  bxwarn/.style={bx, fill=sevone},
  ar/.style={-{Stealth[length=2mm]}, thick},
  grp/.style={draw=black!50, dashed, rounded corners=3pt, inner sep=6pt},
  lbl/.style={font=\scriptsize, text=black!70},
}

% Headings and running heads
\usepackage{titlesec}
\titleformat{\chapter}[display]
  {\normalfont\bfseries\color{figdark}}
  {\large\MakeUppercase{Week}~\thechapter}{4pt}{\LARGE}
\titlespacing*{\chapter}{0pt}{-10pt}{14pt}
\titleformat{\section}{\normalfont\large\bfseries\color{figdark}}{\thesection}{0.6em}{}
\titleformat{\subsection}{\normalfont\normalsize\bfseries}{\thesubsection}{0.6em}{}
\renewcommand{\chaptername}{Week}

\usepackage{fancyhdr}
\pagestyle{fancy}
\fancyhf{}
\fancyheadoffset[R]{5.1cm}      % the header spans the margin column too
\fancyfootoffset[R]{5.1cm}
\fancyhead[L]{\small\itshape\coursename}
\fancyhead[R]{\small\itshape\nouppercase{\leftmark}}
\fancyfoot[C]{\small\thepage}
\renewcommand{\headrulewidth}{0.4pt}
\renewcommand{\chaptermark}[1]{\markboth{Week~\thechapter: #1}{}}
\fancypagestyle{plain}{\fancyhf{}\fancyfoot[C]{\small\thepage}\renewcommand{\headrulewidth}{0pt}}

% Boxes
\usepackage[most]{tcolorbox}
\tcbset{bookbox/.style={enhanced, breakable, boxrule=0.5pt, arc=1.5pt,
  left=6pt, right=6pt, top=4pt, bottom=4pt, fonttitle=\bfseries\small,
  before skip=8pt, after skip=8pt}}
% A definition or a formula to know
\newtcolorbox{definition}[1]{bookbox, colback=figbg, colframe=figdark,
  title={#1}}
% Intuition first: an analogy or the plain-words idea behind a mechanism
\newtcolorbox{intuition}[1][Intuition]{bookbox, colback=white, colframe=figmid,
  title={#1}}
% A worked example; say in the title when it is not from the lecture
\newtcolorbox{worked}[1]{bookbox, colback=white, colframe=figblue,
  title={Worked example: #1}}
% Exam traps and common confusions
\newtcolorbox{traps}[1][Exam traps]{bookbox, colback=sevone, colframe=sevthree,
  title={#1}}
% The end-of-week recap
\newtcolorbox{recap}[1][The week in short]{bookbox, colback=figpale!45,
  colframe=figdark, title={#1}}
% The card at the top of each week: source, lecturer, assignment verdict
\newtcolorbox{lecturecard}{bookbox, colback=black!3, colframe=black!45,
  fontupper=\small}
% Empty-field test: \newcommand macros are \long, so compare with one of them.
\newcommand{\kitempty}{}
\newcommand{\addedtag}{\ (added, not in the lecture)}

% Links, then the glossary (glossaries must come after hyperref)
\usepackage[hidelinks]{hyperref}
\usepackage[toc, nopostdot, nogroupskip]{glossaries}
\makenoidxglossaries
\renewcommand{\glsnamefont}[1]{\textbf{#1}}

% \defterm{key}{Name}{Definition, one or two short sentences.}
% Declare every term in the course's glossary.tex.
\newcommand{\defterm}[3]{\newglossaryentry{#1}{name={#2}, description={#3}}}

% \term{key} or \term[shown text]{key}: the first time a difficult term
% appears in running text. Prints it in bold, puts its definition in the
% margin next to the line and lists the page in the glossary.
% Use it in plain paragraphs only: a margin note cannot leave a box, a table
% or a figure. Inside those, use \gls{key} (no margin note).
\newcommand{\term}[2][]{%
  \textbf{\ifx&#1&\glsentryname{#2}\else#1\fi}\glsadd{#2}%
  \marginpar{\raggedright\footnotesize\textbf{\glsentryname{#2}.}\ \glsentrydesc{#2}}}

% A plain margin remark (a cross-reference, a date, a warning)
\newcommand{\sidenote}[1]{\marginpar{\raggedright\footnotesize\itshape #1}}

% Wide figure: runs into the margin column. Use inside a figure environment.
\newenvironment{wide}{\begin{adjustwidth}{0pt}{-5.1cm}\centering}{\end{adjustwidth}}

% Cross-references printed in bold
\newcommand{\figref}[1]{\textbf{Figure~\ref{#1}}}
\newcommand{\tabref}[1]{\textbf{Table~\ref{#1}}}
\newcommand{\secref}[1]{\S\ref{#1}}
